\scanpage{456}

chooses in each $D'_j$ a ``center'' $a'_j$, in a way compatible with the involution,
then the choice of the $a'_j\in\mathrm{Int}(D'_j)$ defines a torsor structure on
$\partial D'_j$\unsure{\textsuperscript{*}}, and these structures are invariant under
conformal transformations which respect \unsure{the set} of points $a'_j$.

\margin{\uline{N.B.} We no longer assume \unsure{necessarily} that $D_j$ is a circle,
\unsure{only} that $\partial D_j$ not too \deleted{\illegible{1}} \unsure{wild}.
If $\partial D_j$ is a circle, this definition coincides with the preceding one
\unsure{iff} \unsure{$a_j$|$a'_j$} is the center \unsure{of the|of this} circle\ldots}

\uline{Re-N.B.} To define $\tau_j$ on $T_j$, \unsure{considerably}
less data suffice than a topological torsor structure on $T_j$.
Writing \inserted{for $(z,t)\in S_j\times\unsure{\mathbf I|\mathbb I}^{\varepsilon_j}$}
\[
\tau_j(z,t)=(u_t(z),-t)
\tag{37}
\]
where $u_t:S_j\xrightarrow{\sim}S_j$ is a \unsure{homeomorphism} depending continuously on $t$,
writing that $\tau_j|\partial T_j=\tau|\partial T_j$, we find the condition
\[
\text{a)}\quad u_t=\deleted{\unsure{u_1=}}\mathrm{id}
\quad\text{if }t\in\partial\unsure{\mathbf I|\mathbb I}^{\varepsilon_j}\simeq\varepsilon_j
\]

\margin{\unsure{this} is written, if $S_j$ is oriented, $u_1=u_{-1}=\unsure{0}$.}

writing that $\tau_j^2=\mathrm{id}$ we find the condition
\[
\text{b)}\quad u_{-t}=u_t^{-1}
\qquad(t\in\unsure{\mathbf I|\mathbb I}^{\varepsilon_j})
\]
and writing that $T_j^{\tau_j}=\varnothing$ we find
\[
\text{c)}\quad u_0\text{ without fixed points}.
\]

If an orientation is chosen, the $t\longmapsto u_t$ satisfying a) b) c)
correspond to the maps
$[0,1]\xrightarrow{t\mapsto u_t}\deleted{\unsure{\mathrm{Homéom}}}\inserted{\unsure{\mathrm{Aut}}}(\unsure{S_j|S_1})$
such that $u_1=\mathrm{id}$, $u_0$ without fixed point
\deleted{\illegible{2}}\inserted{\uuline{and of order 2}}.
(The case \unsure{considered} above is the one where
$\unsure{t|*}\longmapsto u_{1-t}$ comes from a continuous representation
$\mathbb R\longrightarrow\mathrm{Aut}(\unsure{S_j|S_1})$\ldots)

\scanpage{457}
\mspage{231}

\uline{Remarks} We may set ourselves the task of determining the structure of all anti-involutions \inserted{$\tau$} on an \inserted{oriented} cylinder $T\simeq S\times\unsure{\mathbf I^{\varepsilon}|\mathbb I^{\varepsilon}}$ $(\varepsilon=\operatorname{Or}(S))$ which \deleted{\unsure{induce}} have no fixed point on the boundary --- which already implies that \deleted{they}\inserted{$\tau$} \unsure{permutes|permute} the two connected components of the boundary. More generally, the anti-involutions \inserted{$\tau$} of an oriented \deleted{\unsure{manifold}}\inserted{surface} \uline{with boundary} $\unsure{X}$, having no fixed point on $\partial\unsure{X}$, correspond to manifolds with boundary $\unsure{Y}$ \unsure{equipped} with a subset both \unsure{open} and \unsure{closed} $(\partial Y)'$ of $\partial Y$ --- by associating to such a $\unsure{(Y,(\partial Y)')}$ its \uline{oriented ``\unsure{doubling}''} relative to $(\partial Y)'$ --- to $X,\sigma$ corresponding $(X/\sigma,\operatorname{Im}X^\sigma\longrightarrow X/\sigma)$. If $X$ is connected, \inserted{compact,} $Y$ is \inserted{compact,} \uline{nonorientable}, suppose that \unsure{its type} $(\gamma,j)$, \deleted{let} and \deleted{\unsure{cardinality} $\pi_0$} \unsure{let} \inserted{$0\leq i\leq j$, and,} $\unsure{i|i'}=\operatorname{card}(\pi_0((\partial Y)'))$\inserted{$={}\operatorname{card}\pi_0(X^\sigma)$}. Thus $\chi(Y)=1-(\gamma+j)$, and we \unsure{know|see} at once that \deleted{$\chi(X)=\chi_!(\operatorname{Int}(X))$} (for the $\chi$ of a circle is zero)
\[
\begin{split}
\chi(X)&=\chi_!\left(\underbrace{X\setminus(\unsure{X\mid\partial Y})}_{\unsure{X\mid\operatorname{Int}(Y)}}\right)\\
&=2\chi_!(\operatorname{Int}Y)=2\chi(Y)=2(1-(\gamma+j))
\end{split}
\]
thus if $X$ is of \deleted{genus}\inserted{type} $(g,\nu)$, we will have \inserted{$\chi(X)=2-2g\unsure{-|+}\nu$}
\[
\deleted{\unsure{g+\nu}}\quad 2-2g\unsure{-|+}\nu=2(1-(\gamma+j))
\]
i.e.
\[
(38)\qquad g+\frac{\nu}{2}=\gamma+j
\]

\margin{not true! If $Y$ orientable of type $(\gamma,j)$ (\unsure{same|with} $0\leq i\leq j$) we will have
\[
(39)\qquad\left\{\begin{array}{l}
g+\dfrac{\nu}{2}=2\gamma+j-1\\
\nu=2(j-i)
\end{array}\right.
\]
i.e. $i=j-\dfrac{\nu}{2}$ (as in the case opposite)}

\scanpage{458}

(this requires that $\nu$ be even, which was obvious a priori, since $\sigma$ must permute the elements of $\pi_0(\partial X)$ among themselves without \inserted{there} being a fixed point\ldots)
\deleted{In the present case $g=1,\nu=2$ we therefore \unsure{find|will find}}
\deleted{$\gamma+i=2$}
\deleted{whence the possibilities}
\deleted{$\gamma=2,\ i=0$}
\inserted{But I forgot to note that $\nu$ is determined as a function of $(\gamma,j,i\inserted{={}\unsure{\operatorname{card}((\partial Y)')}})$ by}
\deleted{$\nu=2(j-i)+\unsure{i^{\varepsilon}|i}=2j-i$}
\deleted{$\gamma=2(j-i)$ \unsure{thus}}
\deleted{\unsure{must have}}
\[
(39)\qquad \nu=2(j-i)
\]
i.e. \uline{$i=j-\dfrac{\nu}{2}$}.

Thus here (for $g=0,\nu=2$) we must look for $(\gamma,j,i)$, with $\gamma\in\mathbb N,\ 0\leq i\leq j\in\mathbb N$, such that we have $\nu=2(j-i)$ i.e. $j-i=1$ i.e. $i=j-1$, and $\gamma+j=1$, which gives the \inserted{only} possibility (since $i$ must be $\geq0$ hence $j\geq1$)

a) $\gamma=0,\ j=1,\ i=0$ \unsure{:|.}

\deleted{b) $\unsure{\gamma=1|\gamma=0},\ \unsure{j=0|j=1},\ i=0$.}

Case a) corresponds to the case where \deleted{\illegible{1}} $T^\sigma=\varnothing$, \deleted{it \illegible{2} \unsure{necessarily} a connected component, and} is obtained from the real projective plane by making \deleted{two}\inserted{one} \unsure{hole|holes} in it, \inserted{with boundary (hence Möbius strip)} and by \deleted{bounding one of them by}\inserted{taking the \deleted{\unsure{oriented covering}} oriented doubling}. \deleted{relative to one of these holes.}

The case where $T^\sigma\neq\varnothing$ will necessarily give a quotient \inserted{$Y$} that is \uline{oriented}, we must have for its type $\unsure{(\gamma,j)}$ (since $j\geq1$), the sole case:

\scanpage{459}
\mspage{232}

b) $\gamma=0,\ j=2,\ i=1$ (\uline{oriented} quotient)

obtained from the sphere with two holes i.e. from the cylinder, by taking the oriented doubling \unsure{with respect} to \uline{one} of the holes. \deleted{These are indeed the}

These are indeed the two cases obtained \unsure{resp.\textsuperscript{t}} by $\dot\tau$ and $\dot\tau'$ in formulas (27).

I would now like to construct a situation (under the conditions of the proposition) bringing out a maximum of \unsure{symmetries} --- it is true that we could \unsure{work} with \deleted{\unsure{the} \illegible{1} \unsure{of the}} the \unsure{automorphisms} of $X_0$ \unsure{commuting} with the antipodism\inserted{$\tau$}, \inserted{preserving \unsure{the set} of the $D'_j$,} and respecting (say) topological \unsure{torsors|torsor} structures on \unsure{the set} of the $\partial D'_j$ --- or what \unsure{naturally} amounts to the same (up to an \unsure{indeterminacy} of multiplication by $\unsure{\tau|2}$) the \unsure{automorphisms} of $Y_0$ which \deleted{\illegible{1}} \unsure{preserve|preserving} $D=\bigcup_{i\in I}D_i$, and \unsure{respect|respecting} \unsure{torsor|torsors} structures on \deleted{the boundaries} the connected components $\partial D_i$ of $\partial D$. We expect that (working modulo isotopy) we \unsure{will have|obtain} a group

\margin{no, \uline{without} indeterminacy, by lifting \unsure{defined|to the sphere} so as to \unsure{respect} \unsure{the orientations|the orientation}\ldots}

\scanpage{460}

which will be close to a braid group, and probably computable without much trouble --- and by putting \inserted{\unsure{it}} together with the group generated by the preceding operators, we will perhaps find a start for \unsure{generating} for example \unsure{$\mathcal T_g$|$\Pi_g$} by generators \inserted{(anti-involutive?)} and relations.

Let us consider the \deleted{\unsure{sub}}group $G$ of $A_g=\operatorname{Aut}X_g$ generated by the \deleted{\illegible{1}}\inserted{$\tau_{I'}$} $(I'\in I)$, \deleted{the} Let $\mathcal H$ \deleted{commutators of generators are in} be the subgroup generated by the $\rho_j$ $(j\in I)$. We have
\[
\begin{aligned}
\tau_{I'}\rho_j\tau_{I'}
&=\deleted{\unsure{(\tau_{I'}(\tau_{\{j\}}\tau_{\varnothing})\tau_{I'})}}
\tau_{I'}(\tau_{\{j\}}\tau_{\varnothing})\tau_{I'}\\
&=(\tau_{I'}\tau_{\{j\}})
\underbrace{(\tau_{\varnothing}\tau_{I'})}_{\deleted{\unsure{\rho_{I'_0}}}\unsure{\rho_{I'}}}
\in\mathcal H
\end{aligned}
\]
by formula \unsure{(31)|(32)}, so $\mathcal H$ is an invariant subgroup.

Formulas (32) show that $G/\mathcal H$ is a \uline{commutative} group, \deleted{\unsure{isomorphic to}}\inserted{and even that it is \unsure{isomorphic to}} $\pm1$ by the orientation character (all the $\tau_{I'}$ are equal modulo $\mathcal H$).